%% kapa: exam=1 round=20 year=2009 period=1 form=A subject=civil_law subject_ko=민법 seq=9 number=9
%% source: Q-Net 감정평가사 기출문제 (cst00309) artlSeq=5024178 file=extracted/1008/5024178/2009년 제20회 감정평가사 1차시험 기출문제/1교시.hwp
%% conversion: kordoc 4.24.51 -> tools/gen1.py (auto); review=proofread-sonnet; flags=none
\begin{question}{9}
행위무능력에 관한 설명으로 옳은 것은? (다툼이 있으면 판례에 의함)
\begin{choices}
  \item 미성년자와 거래한 상대방이 미성년자의 취소권을 배제하기 위하여 미성년자의 사술을 주장하는 경우 취소권의 배제를 주장하는 자에게 증명책임이 있다.
  \item 법정대리인이 미성년자에게 특정 재산에 대하여 범위를 정하여 처분을 허락하였다면 그 재산의 처분에 관한 법정대리인의 대리권은 소멸한다.
  \item 행위무능력자와 거래한 상대방이 법률행위 당시 무능력자임을 알고 있었을 경우 상대방은 최고권이나 철회권, 거절권을 행사할 수 없다.
  \item 행위무능력자는 자신의 의사무능력을 이유로 법률행위의 무효를 주장할 수 없다.
  \item 미성년자가 법정대리인의 동의 없이 신용구매계약을 체결한 이후 법정대리인의 동의 없음을 이유로 이를 취소하는 것은 신의칙에 위배된다.
\end{choices}
\end{question}
\answer{1}
